\newpage

% ######################################################################
% PART VI: REFERENCE
% ######################################################################
\part{Reference --- API, Testing, Deployment}

% ######################################################################
% SECTION 29: API REFERENCE
% ######################################################################
\section{Django REST API Reference --- Har Endpoint Ka Pura Detail}

\subsection{Endpoint Map}

\begin{longtable}{p{1.6cm} p{4.2cm} p{7.8cm}}
\toprule
\textbf{Method} & \textbf{Endpoint} & \textbf{Kaam} \\
\midrule
GET & /api/health/ & Server + pipeline readiness check \\
POST & /api/auth/register/ & Naya user + token \\
POST & /api/auth/login/ & Existing user se token \\
POST & /api/research/execute/ & Full research pipeline chalao \\
POST & /api/research/benchmark/ & Single vs multi-agent benchmark \\
GET & /api/benchmark/history/ & Past benchmark evaluations (limit param) \\
\bottomrule
\end{longtable}

\subsection{GET /api/health/}

\begin{codestorybox}
\textbf{Request:} No body. \textbf{Response 200:}
\begin{lstlisting}[style=jsonstyle]
{
  "status": "ok",
  "pipeline": "ready",
  "model": "openai/gpt-oss-120b"
}
\end{lstlisting}

\textbf{Kya karta hai:} Django server + pipeline imports check karta hai. Streamlit app isi se backend reachability decide karta hai (Benchmark Lab ka ``Backend API reachable'' banner).
\end{codestorybox}

\subsection{POST /api/auth/register/}

\begin{codestorybox}
\textbf{Request body:}
\begin{lstlisting}[style=jsonstyle]
{"username": "rahul", "password": "secret123"}
\end{lstlisting}

\textbf{Response 201:}
\begin{lstlisting}[style=jsonstyle]
{"id": 1, "username": "rahul", "token": "abc123..."}
\end{lstlisting}

\textbf{Failures:}
\begin{itemize}
\item 400 --- missing fields ya duplicate username (serializer validation).
\item Password hashed hota hai (\texttt{set\_password}) --- plain text kabhi store nahi hota.
\end{itemize}
\end{codestorybox}

\subsection{POST /api/auth/login/}

\begin{codestorybox}
\textbf{Request body:}
\begin{lstlisting}[style=jsonstyle]
{"username": "rahul", "password": "secret123"}
\end{lstlisting}

\textbf{Response 200:}
\begin{lstlisting}[style=jsonstyle]
{"token": "abc123...", "username": "rahul"}
\end{lstlisting}

\textbf{Failures:} 401 --- galat password ya user exist nahi karta (\texttt{authenticate} check). Response mein 401 deliberately generic hai --- user enumeration prevent karta hai.
\end{codestorybox}

\subsection{POST /api/research/execute/}

\begin{codestorybox}
\textbf{Request body:}
\begin{lstlisting}[style=jsonstyle]
{"topic": "solid state battery market outlook 2026"}
\end{lstlisting}

\textbf{Response 200 (pipeline ke baad):}
\begin{lstlisting}[style=jsonstyle]
{
  "id": 1,
  "topic": "solid state battery market outlook 2026",
  "report": "# Research Report: ... (full markdown)",
  "messages": ["\U0001f50d Researcher found 15 web results...",
               "\U0001f4ca Analyst processed 15 results.",
               "\u2705 Fact-check PASSED (revision 0/3)",
               "\u270d\ufe0f Report complete"],
  "sources_count": 15,
  "revision_count": 0,
  "created_at": "2026-08-17T12:00:00Z"
}
\end{lstlisting}

\textbf{Validation:} Missing topic $\rightarrow$ 400. \textbf{Error path:} Pipeline exception $\rightarrow$ 500, \textbf{no row persisted} (transaction rollback) --- half-saved sessions nahi.
\end{codestorybox}

\subsection{POST /api/research/benchmark/}

\begin{codestorybox}
\textbf{Request body:}
\begin{lstlisting}[style=jsonstyle]
{"topic": "Best practices for RAG in 2026"}
\end{lstlisting}

\textbf{Response 200:}
\begin{lstlisting}[style=jsonstyle]
{
  "id": 1,
  "topic": "Best practices for RAG in 2026",
  "single_agent_report": "...",
  "multi_agent_report": "...",
  "single_agent_depth": 9,
  "multi_agent_depth": 8,
  "single_agent_verifiability": 8,
  "multi_agent_verifiability": 7,
  "verdict": "SINGLE_AGENT_SUPERIOR",
  "created_at": "2026-08-17T12:00:00Z"
}
\end{lstlisting}

\textbf{Verdict logic:} Multi-agent depth+verifiability sum > single-agent sum $\rightarrow$ MULTI\_AGENT\_SUPERIOR; close margin $\rightarrow$ TIE; warna SINGLE\_AGENT\_SUPERIOR.
\end{codestorybox}

\subsection{GET /api/benchmark/history/}

\begin{codestorybox}
\textbf{Query params:} \texttt{?limit=10} (default 20, max 500).

\textbf{Response 200:}
\begin{lstlisting}[style=jsonstyle]
[
  {"id": 1, "topic": "...", "single_agent_depth": 9, "...": "..."},
  {"id": 2, "...": "..."}
]
\end{lstlisting}

\textbf{Ordering:} Newest first (\texttt{-created\_at, -id} tiebreak --- same-timestamp rows deterministic). Empty history $\rightarrow$ \texttt{[]}. Benchmark Lab page isi se charts banata hai.
\end{codestorybox}

\begin{warningbox}
\textbf{Interview note on API design:} (1) Nested resources (\texttt{/api/research/...}) logical grouping ke liye; (2) \texttt{/api/benchmark/history/} alag endpoint isliye kyunki history \textbf{read-only collection} hai --- benchmark POST se alag concern; (3) Response fields consistent snake\_case; (4) Errors: proper status codes (400/401/500) + DRF default error body. Ye REST best practices ke examples hain.
\end{warningbox}

% ######################################################################
% SECTION 30: TESTING DEEP DIVE
% ######################################################################
\section{Testing Deep Dive --- 16 Tests Ka Khoj}

\subsection{Test Infrastructure}

\begin{codestorybox}
\textbf{pytest.ini:}
\begin{lstlisting}[style=bashstyle]
[pytest]
DJANGO_SETTINGS_MODULE = orchestrator_backend.settings
pythonpath = backend
testpaths = backend
\end{lstlisting}

\textbf{Kya karta hai:} (1) Django settings point karta hai; (2) \texttt{backend/} ko sys.path mein daalta hai taaki \texttt{from api.models import ...} chale; (3) sirf \texttt{backend/} mein tests dhoondta hai.
\end{codestorybox}

\subsection{Mocking Strategy --- Real Calls Kabhi Nahin}

\begin{conceptbox}
\textbf{Kyun mock?} Real Tavily/Groq/ChromaDB calls tests ko: slow (seconds-minutes), paid (API credits), non-deterministic (LLM output varies), aur flaky (network) banate hain. Mock se: fast (\textasciitilde9s), free, deterministic, reliable.

\textbf{Kaise mock kiya:}
\begin{itemize}
\item \texttt{mock\_research\_graph} fixture: \texttt{graph.workflow.research\_graph} ko ek stub object se replace --- \texttt{invoke()} canned state return karta hai.
\item \texttt{exploding\_research\_graph}: exception raise karta hai --- 500 error paths test karne ke liye.
\item \texttt{mock\_benchmark}: \texttt{run\_single\_agent\_baseline} + \texttt{evaluate\_outputs} ko stub.
\end{itemize}
\end{conceptbox}

\subsection{Test Case Analysis}

\begin{longtable}{p{2.2cm} p{5.6cm} p{5.8cm}}
\toprule
\textbf{File} & \textbf{Kya test karta hai} & \textbf{Key assertion} \\
\midrule
test\_auth.py (6) & register success/fail, login success/fail & 201+token, 400 dup, 401 wrong pass \\
test\_research.py (4) & execute validation/success/error & 400 no topic, session persisted, 500 no row \\
test\_benchmark.py (6) & benchmark + history + limit & scores persisted, newest first, empty list \\
\bottomrule
\end{longtable}

\begin{interviewbox}
\textbf{Interesting bug jo tests se mila:} History ordering test mein do rows ek hi microsecond mein bane (Django test env time freeze karta hai) --- \texttt{created\_at} tie ho gaya aur ordering nondeterministic. Fix: order by \texttt{-created\_at, -id}. Ye asli robustness win tha --- sirf test ke liye nahi, production ordering ke liye bhi.
\end{interviewbox}

\subsection{How To Run}

\begin{codestorybox}
\begin{lstlisting}[style=bashstyle]
# from project root (venv active)
pytest            # full suite, ~9 seconds

# specific file
pytest backend/api/tests/test_auth.py

# specific test with output
pytest -v backend/api/tests/test_benchmark.py
\end{lstlisting}

\textbf{Interview point:} Test suite \textbf{fast} hone se developers baar baar chalate hain --- 9s suite = daily use. Slow suite (minutes) = kabhi nahi chalta.
\end{codestorybox}

% ######################################################################
% SECTION 31: DEPLOYMENT
% ######################################################################
\section{Deployment --- Demo Se Production Tak}

\subsection{Local Setup (Interview Demo)}

\begin{codestorybox}
\begin{lstlisting}[style=bashstyle]
# 1. Environment
python -m venv .venv
source .venv/bin/activate        # Windows: .venv\Scripts\activate
pip install -r requirements.txt

# 2. Keys (.env ya Streamlit secrets)
GROQ_API_KEY=...
TAVILY_API_KEY=...

# 3. Django backend (port 8001)
cd backend
python manage.py migrate
python manage.py runserver 8001

# 4. Streamlit UI (alag terminal, port 8501)
streamlit run app.py

# 5. MCP server (Claude Desktop ke liye)
python mcp_server.py
\end{lstlisting}

\textbf{Environment separation:} Streamlit secrets (cloud) vs .env (local) --- config/setting.py dono try karta hai.
\end{codestorybox}

\subsection{Streamlit Cloud}

\begin{conceptbox}
\textbf{Steps:} (1) Repo GitHub par push; (2) Streamlit Cloud $\rightarrow$ New app $\rightarrow$ repo select; (3) \textbf{Secrets} mein \texttt{GROQ\_API\_KEY}, \texttt{TAVILY\_API\_KEY} daalo; (4) Deploy. \textbf{Caveat:} Streamlit Cloud ephemeral storage hai --- ChromaDB aur SQLite har session par reset ho sakte hain. Production: managed vector DB + PostgreSQL.
\end{conceptbox}

\subsection{Django Production}

\begin{conceptbox}
\textbf{Production checklist:} (1) \texttt{DEBUG=False}; (2) PostgreSQL; (3) Gunicorn/uvicorn + Nginx reverse proxy; (4) HTTPS (Let's Encrypt); (5) \texttt{ALLOWED\_HOSTS}; (6) CORS allowlist; (7) Rate limiting; (8) Secrets via env, kabhi code mein nahi.
\end{conceptbox}

\subsection{MCP Server for Claude Desktop}

\begin{codestorybox}
\begin{lstlisting}[style=jsonstyle]
// claude_desktop_config.json
{
  "mcpServers": {
    "research-orchestrator": {
      "command": "python",
      "args": ["D:/path/to/mcp_server.py"]
    }
  }
}
\end{lstlisting}

\textbf{Fix yaad rakho:} fastmcp 3.x mein \texttt{FastMCP(description=...)} rename hua \texttt{instructions=...} --- version mismatch par ye error aata hai. (Maine is project mein ye fix kiya.)
\end{codestorybox}

\subsection{Docker (Optional)}

\begin{codestorybox}
\begin{lstlisting}[style=bashstyle]
# Dockerfile (streamlit service)
FROM python:3.11-slim
WORKDIR /app
COPY requirements.txt .
RUN pip install -r requirements.txt
COPY . .
EXPOSE 8501
CMD ["streamlit", "run", "app.py"]
\end{lstlisting}

Compose mein 3 services: \texttt{streamlit} (8501), \texttt{backend} (8001), \texttt{chromadb}. Volumes: ChromaDB persist directory + SQLite.
\end{codestorybox}
